% Standalone sections for inclusion in the combined research paper.
% Requires amsmath, amssymb, amsthm and theorem/lemma/corollary environments.
% Bibliography key: hrw.

\section{A sharp oracle bound for the final Bayesian posterior mean}
\label{sec:alt-final-mean}

A different choice of prior and likelihood gives a substantially faster
approximation guarantee.  The predictor in this section is ordinary
Bayesian model averaging after all observations have been supplied.  It
uses neither a randomly selected training prefix nor an average of the
posteriors at successive times.  The guarantee is in expectation over the
training sample.

The underlying exponential-weights result is closely related to the
recent theorem of H\o gsgaard, Rebeschini and Wegel
\cite{hrw}, who prove optimal expected aggregation
rates for a finite dictionary.  We give a self-contained arbitrary-prior
version, with a conservative temperature constant, so that it applies
directly to continuous neural-network parameter distributions.  The
point of the application below is the explicit neural-network prior and
its dependence on teacher width, depth and weight magnitude.

Let $\pi$ be a probability measure on a measurable parameter space
$\Theta$, and let $(h,x)\mapsto f_h(x)$ be jointly measurable with values
in $[-1,1]$.  Let $D$ be any probability law of $(X,Y)$ with
$Y\in[-1,1]$.  For any bounded measurable predictor $f$, write
$R(f)=\mathbb E_D(f(X)-Y)^2$, and abbreviate $R(h)=R(f_h)$.
Define
\[
 R(h)=\mathbb E_D(f_h(X)-Y)^2,
 \qquad
 \widehat R_S(h)=\frac1n\sum_{i=1}^n(f_h(X_i)-Y_i)^2.
\]
For $n$ independent observations $S\sim D^n$, define
\begin{equation}
 p_S(dh)=
 \frac{\exp\{-\eta n\widehat R_S(h)\}\pi(dh)}
 {\int\exp\{-\eta n\widehat R_S(u)\}\pi(du)},
 \qquad
 \widehat f_S(x)=\int f_h(x)p_S(dh).
 \label{eq:alt-gibbs}
\end{equation}
For $n=0$, the posterior is understood to be $\pi$.

\begin{theorem}[Final posterior mean; arbitrary prior]
\label{thm:alt-final-mean}
Suppose $0<\eta\le \frac14\log(5/4)$.  For every integer $n\ge0$,
\begin{equation}
 \mathbb E_{S\sim D^n}R(\widehat f_S)
 \le
 \inf_Q\left\{
     \int R(h)Q(dh)
       +\frac{\mathrm{KL}(Q\Vert\pi)}{\eta(n+1)}
       \right\}.
 \label{eq:alt-final-oracle}
\end{equation}
The infimum is over probability measures $Q$ independent of $S$; they
may depend on $D$.  In particular, $\eta=1/20$ is admissible.
\end{theorem}

The leading coefficient of the comparison risk in
\eqref{eq:alt-final-oracle} is one.  The theorem allows arbitrary
misspecification and arbitrary dependence of the label on the input.
It concerns the posterior \emph{mean}; the same conclusion is not
asserted for a random posterior draw.  It is also an expectation bound,
not a high-probability bound.

\begin{lemma}[Deleting one observation]
\label{lem:alt-delete-one}
Let $q$ be any probability law on $[-1,1]$, let $y\in[-1,1]$, and set
\[
 p(da)=\frac{e^{-\eta(a-y)^2}q(da)}
                   {\int e^{-\eta(u-y)^2}q(du)}.
\]
If $4(e^{4\eta}-1)\le1$, then
\begin{equation}
 \left(\int a\,q(da)-y\right)^2
       \le\int(a-y)^2p(da).
 \label{eq:alt-delete-one}
\end{equation}
\end{lemma}

\begin{proof}
Write $\ell(a)=(a-y)^2\in[0,4]$, and interpolate from $p$ back to $q$ by
\[
 q_t(da)=\frac{e^{t\eta\ell(a)}p(da)}
                    {\int e^{t\eta\ell(u)}p(du)},
 \qquad 0\le t\le1.
\]
Thus $q_0=p$, $q_1=q$, and boundedness permits differentiation under the
integral:
\[
 \int\ell\,dq-\int\ell\,dp
       =\eta\int_0^1\operatorname{Var}_{q_t}(\ell)\,dt.
\]
The function $a\mapsto(a-y)^2$ is $4$-Lipschitz on $[-1,1]$.
The independent-pair formula for variance gives
$\operatorname{Var}_{q_t}(\ell)\le16\operatorname{Var}_{q_t}(a)$.
Furthermore,
\[
 \frac{dq_t}{dq}\le e^{4\eta(1-t)}.
\]
If probability measures $\nu,\mu$ satisfy $\nu\le K\mu$, then
$\operatorname{Var}_\nu(a)\le K\operatorname{Var}_\mu(a)$: use
$\operatorname{Var}_\nu(a)=\inf_b\int(a-b)^2d\nu$ and choose
$b=\int a\,d\mu$.  Consequently,
\[
 \int\ell\,dq-\int\ell\,dp
 \le16\eta\int_0^1e^{4\eta(1-t)}dt\,
                   \operatorname{Var}_{q}(a)
 =4(e^{4\eta}-1)\operatorname{Var}_{q}(a).
\]
Finally,
$(\int a\,dq-y)^2=\int\ell\,dq-\operatorname{Var}_q(a)$.
The assumed bound on $\eta$ proves the claim.
\end{proof}

\begin{proof}[Proof of Theorem~\ref{thm:alt-final-mean}]
Use $n+1$ independent observations, with full posterior $p$ and posterior
$q_i$ obtained by deleting observation $i$.  Conditional on all these
observations, the image of $p$ under $h\mapsto f_h(X_i)$ is the
exponential tilt in Lemma~\ref{lem:alt-delete-one} of the corresponding
image of $q_i$.  Hence, pointwise in the sample,
\[
 \frac1{n+1}\sum_{i=1}^{n+1}
     \biggl(\int f_h(X_i)q_i(dh)-Y_i\biggr)^2
 \le \int\widehat R_{S_{n+1}}(h)p(dh).
\]
The Gibbs variational principle gives, for every fixed comparison
measure $Q$ with finite relative entropy,
\[
 \int\widehat R_{S_{n+1}}\,dp
 \le \int\widehat R_{S_{n+1}}\,dQ
       +\frac{\mathrm{KL}(Q\Vert\pi)}{\eta(n+1)}.
\]
On taking expectations, each summand on the first display becomes
$\mathbb E_{S\sim D^n}R(\widehat f_S)$: $q_i$ uses $n$ independent
observations and is independent of observation $i$.  The fixed
comparison measure satisfies
$\mathbb E\int\widehat R_{S_{n+1}}dQ=\int R\,dQ$.
This proves the bound for every $Q$, and therefore for its infimum.
For the convenient constant, $e^{1/5}\le(1-1/5)^{-1}=5/4$, so
$\eta=1/20$ satisfies the hypothesis.
\end{proof}

\section{A continuous architecture prior with teacher-dependent rates}
\label{sec:alt-architecture}

We now specify a Bayesian neural-network learner for which the oracle
penalty depends on the teacher's parameters, not the largest architecture
available to the learner.  This is a prior on network architecture and
weights.  It differs from a fixed-width prior that independently draws
all weights from one Gaussian law.

\subsection{Network class and prior}

Inputs lie in any Borel subset $\mathcal X\subseteq[-1,1]^d$, where
$d\ge1$.  Fix a globally Lipschitz activation $\varphi$, and choose
$a\ge\max\{1,\operatorname{Lip}(\varphi)\}$ and
$b=|\varphi(0)|$.  No condition $\varphi(0)=0$ is needed.
A depth-$L$, width-$m$ network, with $L\ge2$ and $m\ge1$, has
$L-1$ hidden layers of width $m$ and a linear readout:
\[
 h^0(x)=(1,x),\qquad h^\ell(x)=\varphi(W^\ell h^{\ell-1}(x)),
 \quad 1\le\ell<L,\qquad z_w(x)=v^{\mathsf T}h^{L-1}(x).
\]
The activation is applied coordinatewise.  The input augmentation
provides first-layer biases; there are no other biases.  The parameter
count is
\begin{equation}
 P_{L,m}=m(d+1)+(L-2)m^2+m.
 \label{eq:alt-parameter-count}
\end{equation}
Its bounded prediction is $f_w(x)=\operatorname{clip}(z_w(x))$, where
$\operatorname{clip}(u)=\max\{-1,\min\{1,u\}\}$.

Let $N\in\mathbb N\cup\{\infty\}$ be the maximum allowed width, and
$L_{\max}\in\{2,3,\ldots\}\cup\{\infty\}$ the maximum allowed depth.  On
allowed pairs $(L,m)$ put prior mass proportional to
\begin{equation}
 \omega_{L,m}=\frac1{L(L-1)m(m+1)},
 \qquad 2\le L\le L_{\max},\quad 1\le m\le N.
 \label{eq:alt-architecture-weight}
\end{equation}
These weights sum to one when $N=L_{\max}=\infty$.  Conditioning on finite
architecture limits only increases the probability of each allowed
pair.  Given $(L,m)$, draw its $P_{L,m}$ weights independently from the
standard Cauchy law, with density $u\mapsto1/[\pi(1+u^2)]$.
Call the resulting proper prior $\pi_{N,L_{\max}}$.
The parameter space is the disjoint union of the allowed architecture
spaces.  In particular, varying $L$ selects networks of genuinely
different depths; it does not require embedding a shallower tanh
network into one fixed depth-$L_{\max}$ architecture.

The learner uses the likelihood
\begin{equation}
 \prod_{i=1}^n\exp\left\{-\frac1{20}
                  (f_w(X_i)-Y_i)^2\right\}
 \label{eq:alt-clipped-likelihood}
\end{equation}
and predicts with the final posterior mean of $f_w$.
This is the usual Gaussian likelihood for a clipped network mean and
observation variance $10$, up to factors independent of $w$.
The theorem below imposes no Gaussian assumption on the actual labels.
In particular it permits arbitrary label noise.  The learner does not
require exact interpolation of the observed labels.

\begin{theorem}[Agnostic neural-network oracle bound]
\label{thm:alt-nn-oracle}
Let $D$ be any probability law on $\mathcal X\times[-1,1]$, and let
$\widehat f_{n,N,L_{\max}}$ be the preceding posterior-mean predictor.  Let
$g=z_{w_*}$ be any network of depth $2\le L\le L_{\max}$ and width
$1\le m\le N$, all of whose weights have absolute value at most
$B\ge1$.  Put
\[
 D_{d,m}=\max\{d+1,m\},\qquad
 A=1+b+aD_{d,m}(B+1),\qquad K=L A^L,
\]
and let $P=P_{L,m}$ as in \eqref{eq:alt-parameter-count}.  For every
$0<\rho\le1$ and $n\ge0$,
\begin{align}
 \mathbb E R(\widehat f_{n,N,L_{\max}})
 &\le R(\operatorname{clip}g)+4\rho \\
 &\quad+\frac{20}{n+1}\left[
  \log\bigl(L(L-1)m(m+1)\bigr)
  +P\log\left(\frac{\pi[1+(B+1)^2]K}{2\rho}\right)
                              \right].
 \label{eq:alt-nn-bound}
\end{align}
The same inequality holds with $R(g)$ in place of
$R(\operatorname{clip}g)$.  Its right side is independent of $N$ and
$L_{\max}$ apart from the requirement that they admit the competitor.
\end{theorem}

The theorem holds for every choice of architecture limits, including
no finite limits at all.  Because it is an expectation bound, it does
not assert a common high-probability sample event for all choices of
these limits.  Taking the infimum in \eqref{eq:alt-nn-bound} permits a
different competitor at every sample size.

\begin{lemma}[Parameter perturbations]
\label{lem:alt-perturb}
In the setting of Theorem~\ref{thm:alt-nn-oracle}, if
$\|w-w_*\|_\infty\le r\le1$, then
\[
 \sup_{x\in\mathcal X}|z_w(x)-z_{w_*}(x)|\le K r.
\]
\end{lemma}

\begin{proof}
Both parameter vectors have coordinates bounded by $B+1$.
Let $M_\ell$ bound the maximum coordinate magnitude in hidden layer
$\ell$ of either network, and let $\delta_\ell$ bound their coordinatewise
difference, uniformly in $x$.  Since $M_0\le1$, the activation bound
$|\varphi(u)|\le b+a|u|$ gives $M_\ell\le A^\ell$ by induction.
The perturbation recurrence is
\[
 \delta_\ell\le aD_{d,m}r A^{\ell-1}
                    +aD_{d,m}(B+1)\delta_{\ell-1},
 \qquad \delta_0=0.
\]
Both $aD_{d,m}$ and $aD_{d,m}(B+1)$ are at most $A$, so induction gives
$\delta_\ell\le \ell r A^\ell$.
For the readout, use $D_{d,m}\le A$ and $D_{d,m}(B+1)\le A$ to obtain
\[
 |z_w-z_{w_*}|
 \le D_{d,m}r A^{L-1}+D_{d,m}(B+1)(L-1)r A^{L-1}
 \le LrA^L.
\]
\end{proof}

\begin{proof}[Proof of Theorem~\ref{thm:alt-nn-oracle}]
Choose a comparison measure supported on architecture $(L,m)$, uniform
on the parameter cube
$\{w:\|w-w_*\|_\infty\le r\}$ with $r=\rho/K\le1$.
Lemma~\ref{lem:alt-perturb} and the fact that clipping is nonexpansive
show $\|f_w-\operatorname{clip}g\|_\infty\le\rho$ throughout this cube.
Squared loss is $4$-Lipschitz in its prediction on $[-1,1]$, hence
\[
 \int R(f_w)Q(dw)\le R(\operatorname{clip}g)+4\rho.
\]
On the cube, each Cauchy coordinate density is at least
$1/\{\pi[1+(B+1)^2]\}$.  The uniform density there is $(2r)^{-P}$.
Moreover, the architecture has prior mass at least $\omega_{L,m}$.
Therefore
\[
 \mathrm{KL}(Q\Vert\pi_{N,L_{\max}})
 \le \log\frac1{\omega_{L,m}}
       +P\log\left(\frac{\pi[1+(B+1)^2]}{2r}\right).
\]
Apply Theorem~\ref{thm:alt-final-mean} with $\eta=1/20$ and substitute
$r=\rho/K$.  Finally, clipping cannot increase squared error against a
label in $[-1,1]$, so $R(\operatorname{clip}g)\le R(g)$.
\end{proof}

\subsection{The same rate with conditionally Gaussian weights}

The logarithmic cost of the teacher's weight magnitude does not require
Cauchy coordinates.  There is also a Gaussian scale-mixture version of
the prior.  After selecting $(L,m)$ with the same architecture law,
choose an integer $j\ge0$ with probability $1/[(j+1)(j+2)]$, and then
draw the weights independently from $N(0,4^j)$.  The prior is proper,
and its coordinates are independent conditional on the architecture
and scale.

\begin{corollary}[Gaussian scale mixture]
\label{cor:alt-gaussian-mixture}
Use this Gaussian scale-mixture prior with the likelihood and predictor
of \eqref{eq:alt-clipped-likelihood}.  For the competitor in
Theorem~\ref{thm:alt-nn-oracle}, let
$j_B=\lceil\log_2 B\rceil$.  Then
\begin{align}
 \mathbb E R(\widehat f_{n,N,L_{\max}})
 &\le R(\operatorname{clip}g)+4\rho\\
 &\quad+\frac{20}{n+1}\left[
   \log(1/\omega_{L,m})
   +\log\bigl((j_B+1)(j_B+2)\bigr)
   +P\log\left(\frac{B\sqrt{2\pi}e^2K}{\rho}\right)
                           \right].
 \label{eq:alt-gaussian-mixture}
\end{align}
\end{corollary}

\begin{proof}
The comparison measure selects architecture $(L,m)$, selects scale
$j_B$, and uses the same uniform parameter cube of radius $r=\rho/K$.
Put $s=2^{j_B}$, so $B\le s\le2B$.
Every coordinate in the cube satisfies $|w_i|\le B+1\le2s$; its
$N(0,s^2)$ density is therefore at least
$e^{-2}/(s\sqrt{2\pi})$.  The coordinate contribution to relative
entropy is at most
\[
 P\log\left(\frac{s\sqrt{2\pi}e^2}{2r}\right)
 \le P\log\left(\frac{B\sqrt{2\pi}e^2}{r}\right).
\]
Adding the architecture and scale selection costs gives the bracket
in \eqref{eq:alt-gaussian-mixture}.  The risk comparison and final
posterior-mean theorem apply as before.
\end{proof}

The additional scale cost is only
$O(1+\log(1+\log B))$.  Both prior choices consequently have the
same polynomial width--depth rates described next.  A single fixed
Gaussian scale would instead charge order $PB^2$ for translating the
comparison cube; the scale mixture makes this cost logarithmic in $B$.

\subsection{Rates, activations, and what changes in the learner}

Setting $\rho=1/(n+1)$ in \eqref{eq:alt-nn-bound} yields a comparison
penalty of order
\begin{equation}
 \frac{1}{n+1}\left\{
 1+\log(Lm)
 +P_{L,m}\left[
   \log(n+1)+\log(B+1)+\log L
      +L\log A
                  \right]\right\}.
 \label{eq:alt-rate-summary}
\end{equation}
Here $A=1+b+a\max\{d+1,m\}(B+1)$, as above, and the hidden
constant is numerical.  Thus the cost is approximately one
bit per parameter per required precision bit, together with the
architecture description.  The dependence on the weight bound is
logarithmic.  No bound on $B$ must be selected in advance by the learner.

For fixed activation, input dimension and depth, competitors of width
$m\le n^\alpha$ and weights bounded by $n^\beta$, where $\beta$ is any
fixed nonnegative constant, have excess expected risk
$O(n^{2\alpha-1}\log n)$.  Every $\alpha<1/2$ is therefore admissible.
For a single hidden layer ($L=2$),
$P_{2,m}=m(d+2)$, and this improves to
$O(n^{\alpha-1}\log n)$: every $\alpha<1$ is admissible.
If the depth also varies, the displayed bound has polynomial, rather
than exponential, depth dependence.  Suppressing logarithms and holding
$d$, $B$ and the activation fixed, it is
$\widetilde O((mdL+m^2L^2)/(n+1))$.
These statements are approximation guarantees: the target distribution
need not be representable by any network.

The activation assumption includes tanh and ReLU with $a=1,b=0$;
GELU, $\varphi(u)=u\Phi(u)$, with $a=2,b=0$; and softplus with
$a=1,b=\log2$.  In particular, the proof neither relies on bounded
hidden activations nor requires $\varphi(0)=0$.  The Gaussian derivative
bound $|\Phi(u)+u\phi(u)|\le1+(2\pi e)^{-1/2}<2$ verifies the stated
GELU constant.  Internal biases can be added by augmenting each layer
with a constant coordinate and changing the parameter count and
$D_{d,m}$ accordingly.

The price of this improved bound is a change of prior and prediction
convention, not an unproved property of the original dense Gaussian
prior.  A draw from $\pi_{N,L_{\max}}$ first selects the \emph{active}
architecture and then draws its weights.  At a fixed depth and maximum width it
can equivalently be regarded as a structural sparsity prior, with exact
zero connections outside the selected subnetwork.  It is not a product
law on all the weights of a dense width-$N$ network.  Its preference for
small architectures is explicit.  The guarantee therefore does not
resolve whether an unrestricted dense Gaussian prior has a comparable
polynomial-depth approximation theorem.

Finally, the loss in the likelihood is computed from the clipped
network output, and the predictor averages those clipped outputs.
Replacing either with an unrestricted real output is a different
learner.  Boundedness is used in Lemma~\ref{lem:alt-delete-one}; no such
replacement follows from that lemma.
